%%%
%
% $Autor: Müller, Hinrichs $
% $Datum: 2026-06-09 14:18:35Z $
% $Pfad: Radar.tex
% $Version: 1 $
%
% Project: 
%
%%%


\chapter{\textcolor{red}{Python Tool XY}}

\section{Mikrocontroller-Programmierung mit der Arduino IDE}

\subsection{Allgemeines zur Arduino IDE}
Die Arduino Integrated Development Environment (IDE) ist eine plattformunabhängige Open-Source-Entwicklungsumgebung, die speziell für die Programmierung von Mikrocontrollern entworfen wurde. Sie basiert im Kern auf der Programmiersprache C/C++, verbirgt jedoch viele komplexe hardwarenahe Konfigurationen hinter einer stark vereinfachten Syntax und einer umfangreichen Standardbibliothek. 

Ein klassisches Arduino-Programm (auch ,,Sketch`` genannt) ist strukturell immer in zwei wesentliche Hauptfunktionen unterteilt: Die \texttt{setup()}-Funktion, welche genau einmal beim Systemstart ausgeführt wird und für die Initialisierung von Pins und Schnittstellen zuständig ist, sowie die \texttt{loop()}-Funktion, welche im Anschluss als Endlosschleife die eigentliche Arbeitslogik des Mikrocontrollers abarbeitet. Die fertigen Programme werden direkt aus der IDE heraus kompiliert und über eine serielle USB-Verbindung auf den Flash-Speicher des Mikrocontrollers übertragen.

\subsection{Das Steuerungsprogramm des Ultraschallradars}
Der Quellcode für den Arduino Nano 33 BLE übernimmt die direkte Hardwaresteuerung der Sensorik, der Aktuatorik sowie der visuellen Statusanzeigen (LEDs). Bevor das System in den regulären Betriebsmodus wechselt, durchläuft es in der \texttt{setup()}-Phase einen sogenannten Power-On Self-Test (POST). Hierbei teilt der Mikrocontroller der verbundenen Processing-Software zunächst den Status \texttt{S:TESTING} mit. Daraufhin werden die LEDs auf Funktion geprüft, der Servomotor absolviert eine Testdrehung und der Ultraschallsensor führt eine initiale Referenzmessung durch. 

Nur wenn diese Testmessung valide Werte (zwischen 1\,cm und 400\,cm) liefert, wird das System freigegeben, die grüne Status-LED aktiviert und das Signal \texttt{S:OK} an den Computer gesendet. Tritt ein Fehler auf, beispielsweise durch einen Verbindungsabbruch zum Sensor (Timeout), leuchtet die rote LED auf, das System blockiert aus Sicherheitsgründen den weiteren Betrieb und meldet \texttt{S:ERR\_SENSOR}.

Im laufenden und fehlerfreien Betrieb (\texttt{loop}) iteriert das System kontinuierlich durch die Positionen des Servomotors (0\textdegree{} bis 180\textdegree{} und zurück). Nach jedem Positionsschritt pausiert das System kurz, um die mechanische Trägheit des Motors auszugleichen. Anschließend feuert der HC-SR04-Sensor einen 10\,µs langen Ultraschall-Impuls ab und misst die Zeit bis zum Eintreffen des Echos. Aus dieser Laufzeit errechnet der Mikrocontroller unter Einbeziehung der Schallgeschwindigkeit die genaue Distanz. 

Abschließend formatiert der Arduino diese Wertepaare in das speziell für dieses Projekt definierte Kommunikationsprotokoll (\texttt{D:Winkel,Distanz}) und übermittelt sie in Echtzeit an die serielle Schnittstelle zur weiteren grafischen Auswertung.

\newpage
\subsection{Quellcode der Arduino-Steuerung}

% HINWEIS ZUM PFAD: 
% Passe diesen Pfad genau wie bei der Processing-Datei an deine Ordnerstruktur an!
% Beispiel: \lstinputlisting[...]{../../Code/Arduino/RadarControl.ino}

\lstinputlisting[
language=C++, 
caption={Arduino Steuerungscode mit Doxygen-Dokumentation}, 
label={lst:arduino_radar}, 
numbers=left, 
frame=single, 
basicstyle=\footnotesize\ttfamily, 
keywordstyle=\color{blue}, 
commentstyle=\color{teal}, 
stringstyle=\color{red}
]{../../Code/Ultraschallradar/UltraschallradarArduinoIDEMitDoxygen/UltraschallradarArduinoIDEMitDoxygen.ino}